\documentclass[11pt]{article}

\usepackage[T1]{fontenc}
\usepackage[utf8]{inputenc}
\usepackage{lmodern}
\usepackage[french]{babel}
\usepackage[a4paper,margin=3cm]{geometry}
\usepackage{microtype}
\usepackage{parskip}

\pagestyle{empty}

\begin{document}

\begin{center}
{\LARGE Une porte de village qui apprend à faire attendre}\\[0.6em]
{\large Explication accessible du « contrôle d'admission stigmergique et
  de quorum »}\\[0.3em]
{\small \texttt{moderation-pipeline-prototype} --- août 2026}
\end{center}

\vspace{1em}

\textbf{Le problème.} Certains litiges qui arrivent devant un système de
modération sont réellement difficiles à trancher~: chaque camp avance des
arguments qui attaquent ceux de l'autre, et la méthode rapide et peu
coûteuse censée les départager ne parvient à aucune conclusion nette.
Dans ce cas, le système se rabat sur une méthode bien plus minutieuse,
et bien plus coûteuse, pour démêler l'ensemble --- une méthode utile,
mais que l'on ne souhaite pas offrir gratuitement, sur simple demande, à
quiconque la sollicite. Le défaut se situe dans la manière dont cette
méthode coûteuse se déclenche. Dans la conception d'origine, la toute
première personne à se présenter avec un argument d'apparence plausible
fait basculer le dossier vers le traitement coûteux, immédiatement.
C'est exactement le genre de règle qu'un attaquant adore~: se présenter
une seule fois, en masse, avec un tas d'arguments fabriqués, et la porte
coûteuse s'ouvre gratuitement, autant de fois qu'on le souhaite.

\textbf{L'idée.} Le remède emprunte deux mécanismes à la nature, et il
vaut la peine de les prendre l'un après l'autre, car chacun résout une
moitié différente du problème.

Le premier mécanisme est la façon dont une colonie de fourmis retient
quel chemin vers la nourriture mérite d'être suivi. Une fourmi ne vote
pas une fois pour toutes~; elle dépose une odeur en marchant, et cette
odeur s'estompe régulièrement, à moins que d'autres fourmis, arrivées
indépendamment, n'empruntent le même chemin et ne la raniment. Un chemin
ne reste marqué que s'il continue de gagner des passages frais, venus de
fourmis différentes --- une seule fourmi qui fait des allers-retours sur
sa propre piste toute la journée laisse, en pratique, une trace tout
aussi faible qu'un unique passage. Transposé ici~: chaque dossier
disputé porte un score qui s'affaiblit un peu à chaque pas de l'horloge,
et un nouvel argument ne relève ce score que d'un tout petit peu s'il
provient d'une source ayant déjà pris la parole récemment. Une
préoccupation véritablement nouvelle, indépendante, maintient un dossier
en vie~; le même plaignant répétant la même chose en boucle, non.

Le second mécanisme est la façon dont les bactéries décident, en
colonie, du moment où elles sont assez nombreuses pour agir ensemble ---
un mécanisme appelé détection de quorum. Aucune bactérie isolée ne prend
cette décision seule~; toute la colonie lit un signal chimique partagé
qui reflète la densité déjà atteinte, et relève ou abaisse en
conséquence son seuil collectif d'action. Transposé ici~: la barre qu'un
dossier doit franchir pour atteindre le traitement coûteux n'est pas
fixe. Elle s'élève automatiquement dès que ce traitement coûteux est
déjà occupé par d'autres dossiers --- exactement comme une salle de
restaurant devient plus difficile à intégrer à mesure qu'elle se
remplit, sans que personne à l'entrée ne décide consciemment de durcir
les critères.

\textbf{L'assemblage.} Le score d'un dossier disputé s'estompe comme une
odeur de piste qui s'évente, n'est relevé que par un renforcement
d'apparence nouvelle, est plafonné pour qu'aucun instant isolé
d'activité ne puisse le faire bondir au-delà d'une limite raisonnable, et
doit franchir une barre qui monte avec l'occupation déjà en cours du
traitement coûteux. Une seule salve d'indignation fabriquée, aussi
massive soit-elle, ne peut faire franchir à un score plafonné et
décroissant un seuil en un seul instant --- il n'y a tout simplement pas
assez de poids autorisé à entrer dans un seul pas de l'horloge, quel que
soit le nombre de voix factices qui se présentent d'un coup.

\textbf{Ce qui a réellement été mesuré, non simplement espéré.} Le
mécanisme a été confronté à son prédécesseur --- la règle d'admission
instantanée --- sous plusieurs formes d'attaque, et la comparaison a été
menée honnêtement, y compris sur le cas où le nouveau mécanisme ne
l'emporte pas complètement. Une salve unique, de taille quelconque, d'une
seule voix à cinq mille, arrivant en un instant~: entièrement bloquée, à
chaque fois, là où l'ancienne règle laissait passer chacune d'entre
elles. La même salve soutenue très longtemps depuis une seule source
repérable~: toujours entièrement bloquée, et de façon démontrable, non
pas seulement « bloquée jusqu'ici dans les essais menés ». Deux cents
litiges distincts, chacun frappé d'une salve isolée, aux côtés de quinze
litiges authentiques en croissance organique au même moment~: l'ancienne
règle laissait passer toutes les attaques et admettait aussi tous les
litiges authentiques, pour un coût de traitement réel de 215 dossiers~;
le nouveau mécanisme n'en laissait passer aucune tout en admettant les
quinze litiges authentiques, pour un coût réel de 15 dossiers --- quatre-vingt-treize
pour cent de travail en moins imposé au traitement coûteux. L'exception
honnête, énoncée sans détour plutôt que dissimulée~: un attaquant patient,
disposé à fabriquer une identité fraîche, jamais réutilisée et
d'apparence indépendante à chaque pas de l'horloge, pendant plusieurs pas
d'affilée, n'est pas complètement arrêté. Cet attaquant finit par passer
tout de même, au prix d'un effort et d'un temps environ cinq fois
supérieurs à ce que l'ancienne règle lui aurait coûté. Le mécanisme
transforme une attaque gratuite en attaque coûteuse~; il ne la rend pas
impossible, et cette limite est rapportée ici plutôt que laissée à la
découverte, plus tard, par quelqu'un d'autre.

\textbf{Ce que ce mécanisme ne fait délibérément pas.} Il ne dit rien de
qui a \emph{raison} dans le litige lui-même --- ce jugement reste
entièrement confié au traitement coûteux qu'il protège, une fois le
dossier admis. Il suppose qu'il existe un coût réel, même modeste, à se
présenter comme une voix toute nouvelle et indépendante à chaque pas~;
si cette hypothèse cesse de tenir pour un attaquant donné, le retard
imposé constitue toute la défense, non un filet posé sur quelque chose
de plus solide. Et le tour de l'estompage plafonné d'une part, celui de
la barre qui monte d'autre part, pris chacun tout à fait à part, sont
individuellement bien connus en dehors de la biologie aussi --- le
premier ressemble à un limiteur de débit ordinaire, utilisé partout en
informatique~; le second ressemble à une technique de réseau vieille de
plusieurs décennies, destinée à désengorger une file avant qu'elle ne
déborde. Ce qu'une recherche approfondie dans la littérature de recherche
pertinente n'a pas fait apparaître, c'est cette combinaison précise des
trois éléments, gardant spécifiquement ce type de recours coûteux et
difficile à paralléliser. C'est là l'affirmation restreinte qui est
faite, et seulement dans la mesure où cette recherche a été menée à
fond --- non une prétention que personne, nulle part, n'aurait jamais
combiné ces idées auparavant.

\textbf{Pourquoi cela mérite d'être consigné.} Rien, dans le traitement
coûteux lui-même, n'a eu besoin de changer~; le correctif se loge
entièrement à la porte qui le précède. Les ingrédients pris
individuellement --- un score qui s'estompe, un plafond sur ce qu'un
seul instant peut apporter, une barre qui monte avec l'affluence --- sont
chacun empruntés à des domaines forts de leur propre longue expérience,
biologique aussi bien que informatique. Ce qui est rapporté ici, c'est la
manière dont ils se composent, testée honnêtement contre l'attaque
qu'ils devaient contrer, avec le cas où le succès n'est que partiel
énoncé aussi clairement que les cas de succès complet.

\end{document}
